\begin{table*}[t]
  \centering
  \caption{不同无线边缘场景下的方法选择}
  \label{tab:scenario-selection}
  \begin{threeparttable}
    \small
    \setlength{\tabcolsep}{4pt}
    \renewcommand{\arraystretch}{1.28}
    \begin{tabularx}{\textwidth}{
      @{}>{\centering\arraybackslash}p{2.05cm}
      >{\raggedright\arraybackslash}p{2.45cm}
      >{\centering\arraybackslash}p{2.65cm}
      >{\raggedright\arraybackslash}X
      >{\raggedright\arraybackslash}X@{}}
      \toprule
      主导瓶颈 & 典型场景 & 优先考虑的方法 & 选择理由 & 使用时需要注意的问题 \\
      \midrule
      同步等待与终端异构
        & 大规模物联网；设备能力和在线状态差异明显
        & 客户端选择＋资源分配
        & 同时控制“谁参与”和“如何完成本轮传输”，缓解慢终端造成的等待
        & 快设备优先会造成选择偏差；状态采集和联合求解也有开销\cite{Xu2021ClientBandwidth,Nishio2019FedCS} \\
      高维更新与窄上行
        & 视觉模型、移动终端或低速上行链路
        & 量化、稀疏化及误差反馈
        & 直接减少每轮需要编码和发送的更新量
        & 需计入索引、码本和残差；压缩误差可能增加训练轮数\cite{Reisizadeh2020FedPAQ,Sattler2020STC} \\
      密集终端同时接入
        & 同步条件可控、聚合函数可加的蜂窝或传感网络
        & AirComp；必要时结合稀疏化
        & 在同一时频资源中完成多用户叠加，避免逐用户正交接入
        & 依赖信道估计、同步与功率可行性，并引入聚合误差\cite{Zhu2020BroadbandAnalog,Oh2024SparseOneBitOTA} \\
      回传链路受限
        & 多小区、工业园区或云—边—端集群
        & 分层联邦＋两级压缩
        & 先减少边—云同步频率，再降低簇模型的回传载荷
        & 需联合控制量化误差、簇间漂移和边缘负载\cite{Liu2023HierarchicalQuantization,Zhong2022HierarchicalHeterogeneous} \\
      多瓶颈随时间变化
        & 移动、多业务竞争或时变信道环境
        & 跨轮选择、压缩和资源联合优化
        & 可随训练阶段和网络状态切换主要优化变量
        & 在线决策、反馈和模型迁移可能抵消部分通信收益\cite{Xu2021ClientBandwidth,Marnissi2024AdaptiveSQ,Azimi2024HierarchicalOTA} \\
      \bottomrule
    \end{tabularx}
  \end{threeparttable}
\end{table*}
